% !TeX root = ../main.tex

% ==================================================
% A.10. 여러 데코레이터의 평가·적용 순서
% ==================================================

\subsection{여러 데코레이터의 평가·적용 순서}

Decorator Stacking · Evaluation Order · Application Order · Wrapper Call Chain · Function Composition

\subsubsection{이전 학습과의 연결}

\zsh{A.9}에서는 인수를 받는 데코레이터를 다음 세 단계로 분석했다.

\begin{dhkcode}
@repeat(3)
def greet():
    print("Hello")
\end{dhkcode}

\begin{dhksmry}
1. repeat(3)을 호출해 데코레이터 함수 객체를 생성한다.

2. 반환된 데코레이터에 원본 greet 함수 객체를 전달해 wrapper를 생성한다.

3. 이후 greet()을 호출하면 wrapper가 실행된다.
\end{dhksmry}

이 구조를 이해했다면 이제 여러 데코레이터를 하나의 함수에 적용할 수 있다.

\begin{dhkcode}
@first
@second
def greet():
    print("Hello")
\end{dhkcode}

이 코드는 단일 데코레이터의 변환식을 확장한 것으로, 개념적으로 다음과 대응한다.

\begin{dhkcode}
greet = first(second(greet))
\end{dhkcode}

그런데 여기서 주의할 점이 있다.

\zsh{first}가 위에 작성되어 있으므로 먼저 적용될 것이라고 생각하기 쉽지만, 실제로는 그렇지 않다.

Python에서는 데코레이터 표현식의 평가 순서, 함수 객체에 데코레이터를 적용하는 순서, 최종 함수 호출 시 wrapper가 실행되는 순서를 구별해야 한다.

이 세 순서는 서로 관련되어 있지만 동일하지 않다.

\subsubsection{이번 학습의 핵심 질문}

이번 학습에서는 다음 문제를 해결할 것이다.

\begin{dhksmry}
1. 여러 데코레이터 표현식은 어떤 순서로 평가되는가?

2. 평가된 데코레이터는 원본 함수 객체에 어떤 순서로 적용되는가?

3. 데코레이션 이후 함수 이름은 어떤 wrapper 객체를 참조하는가?

4. 실제 함수를 호출하면 ㅈrapper는 어떤 순서로 실행되는가?

5. 반환값과 예외는 wrapper 사이에서 어떤 경로로 전달되는가?

6. __wrapped__와 __closure__를 이용하면 중첩 구조를 어떻게
   추적할 수 있는가?

7. 데코레이터의 적용 순서를 바꾸면 프로그램의 동작이 달라질 수 있는가?
\end{dhksmry}

핵심 원리는 다음과 같다.

\begin{dhksmry}[centered]
데코레이터 표현식은 위에서 아래로 평가되고, 

그 결과는 아래에서 위로 적용된다.

이후 실제 호출은 가장 바깥쪽 wrapper부터 안쪽으로 진행하며,

정상 반환은 반대 방향으로 전달된다.
\end{dhksmry}

다만 모든 데코레이터가 반드시 wrapper를 반환하는 것은 아니므로, 마지막의 호출·반환 순서는 함수 기반 wrapper를 중첩한 경우에 성립하는 구조다.

\subsubsection{세 가지 순서를 먼저 구별하자}

다음 코드를 기준으로 생각해 보자.

\begin{dhkcode}
@A()
@B()
def func():
    ...
\end{dhkcode}

\zsh{A()}와 \zsh{B()}가 각각 데코레이터 함수 객체를 반환한다고 가정한다.

그러면 세 단계는 다음과 같이 진행 된다.

\begin{dhknote}
\begin{dhksmry}[centered]
1. 데코레이터 표현식 평가 - 위에서 아래
\end{dhksmry}
\end{dhknote}

\begin{dhksmry}
    (1) A() : 데코레이터 A 생성

    (2) B() : 데코레이터 B 생성
\end{dhksmry}

\begin{dhknote}
\begin{dhksmry}[centered]
2. 원본 함수에 적용 - 아래에서 위
\end{dhksmry}
\end{dhknote}

\begin{dhksmry}[centered]
original \ra B(original) \ra A(B(original))
\end{dhksmry}

\begin{dhknote}
\begin{dhksmry}[centered]
3. 실제 호출 - 바깥에서 안으로, 반환은 반대로
\end{dhksmry}
\end{dhknote}

\begin{dhkgraph}[uniform]
node n1 (1,1) :: 원본 함수
group g1(n1) :: B wrapper
group g2(g1) :: A wrapper

explain :: 호출 진입은 A wrapper에서 원본 함수순으로,//반환 전달은 원본 함수에서 A wrapper 순으로

caption :: 원본 함수에 가까운 wrapper가 먼저 생성되고, 마지막에 생성된 바깥쪽 wrapper가 최종 함수 이름에 바인딩된다.
\end{dhkgraph}

이를 표로 정리하면 다음과 같다.

\begin{dhktable}
|구분	|순서	|의미|
|:---|:---|:---|
|표현식 평가|	A() → B()|	데코레이터 객체 준비|
|데코레이터 적용|	B → A|	원본 함수에서 바깥쪽으로 객체 변환|
|실제 호출 진입|	A wrapper → B wrapper → 원본|	바깥쪽에서 안쪽으로|
|정상 반환 전달|	원본 → B wrapper → A wrapper|	안쪽에서 바깥쪽으로|
\end{dhktable}

여기서 첫 번째와 두 번째 순서가 반대라는 점이 중요하다.

이것을 단순히 "데코레이터는 역순으로 실행된다"라고 암기하면 안 된다.

무엇을 평가하는지, 무엇을 적용하는지, 어떤 함수 객체를 실제로 호출하는지에 따라 순서가 달라지기 때문이다.

\subsubsection{실행 결과로 세 가지 순서를 검증하기}

먼저 하나의 독립적인 예제로 평가·적용·호출 순서를 모두 확인해 본다.

\begin{dhkcode}
from functools import wraps

def layer(label):
    print("evaluate:", label)

    def decorator(func):
        print("apply:", label)

        @wraps(func)
        def wrapper(*args, **kwargs):
            print("enter:", label)

            result = func(*args, **kwargs)

            print("exit:", label)
            return result

        return wrapper

    return decorator


@layer("A")
@layer("B")
def calculate(x):
    print("body:", x)
    return x * 2


print("definition complete")
print("result:", calculate(3))
\end{dhkcode}

예상 출력:

\begin{dhkzsh}
evaluate: A
evaluate: B
apply: B
apply: A
definition complete
enter: A
enter: B
body: 3
exit: B
exit: A
result: 6
\end{dhkzsh}

이 코드는 세 가지 시점으로 나누어 분석해야 한다.

\begin{dhknote}
\begin{dhksmry}[centered]
1. 데코레이터 표현식 평가
\end{dhksmry}
\end{dhknote}

Python이 함수 정의문을 실행하면서 먼저 다음 표현식을 평가한다.

\begin{dhkcode}
layer("A")
\end{dhkcode}

\zsh{layer}의 매개변수 \zsh{label}이 문자열 객체 \zsh{"A"}에 바인딩된다.

\zsh{layer}는 중첩된 \zsh{decorator} 함수 객체를 생성하고 반환한다.

이때 다음 출력이 발생한다.

\begin{dhkzsh}
evaluate: A
\end{dhkzsh}

이어서 두 번째 데코레이터 표현식을 평가한다.

\begin{dhkcode}
layer("B")
\end{dhkcode}

마찬가지로 \zsh{"B"}를 참조하는 클로저 환경을 가진 새로운 \zsh{decorator} 함수 객체가 생성된다.

\begin{dhkzsh}
evaluate: B
\end{dhkzsh}

여기서 두 데코레이터 함수 객체는 같은 \zsh{decorator} 정의문으로부터 생성되었지만 서로 다른 함수 객체다.

각각 다른 \zsh{label} 바인딩을 유지하고 있다.

\begin{dhknote}
\begin{dhksmry}[centered]
2. 원본 함수 객체에 데코레이터 적용
\end{dhksmry}
\end{dhknote}

데코레이터 표현식의 평가가 끝나면 원본 \zsh{calculate} 함수 객체가 생성된다.

이제 아래쪽 데코레이터부터 적용된다.

먼저 \zsh{layer("B")}가 반환한 데코레이터에 원본 함수 객체가 전달된다.

\begin{dhkcode}
inner_wrapper = decorator_B(original_calculate)
\end{dhkcode}

이때 다음 출력이 발생한다.

\begin{dhkzsh}
apply: B
\end{dhkzsh}

\zsh{decorator_B}는 새로운 wrapper 함수 객체를 생성하고 반환한다.

이 wrapper는 \zsh{label="B"}와 원본 \zsh{calculate} 함수 객체에 접근할 수 있다.

그다음 위쪽 데코레이터가 적용된다.

\begin{dhkcode}
outer_wrapper = decorator_A(inner_wrapper)
\end{dhkcode}

여기서 \zsh{decorator_A}가 전달받는 \zsh{func}는 

원본 \zsh{calculate}가 아니라 B의 wrapper 함수 객체다.

이때 다음 출력이 발생한다.

\begin{dhkzsh}
apply: A
\end{dhkzsh}

마지막으로 \zsh{outer_wrapper}가 모듈 이름 \zsh{calculate}에 바인딩된다.

따라서 최종 관계는 다음과 같다.

\begin{dhkbranch}[in=8]
node main :: calculate
arrow main
node main :: A의 wrapper
line main
line main :: func
arrow main
node main :: B의 wrapper
line main
line main :: func
arrow main
node main :: 원본 calculate
\end{dhkbranch}

\begin{dhknote}
\begin{dhksmry}[centered]
3. 실제 함수 호출
\end{dhksmry}
\end{dhknote}

이제 다음 문장이 실행된다.

\begin{dhkcode}
calculate(3)
\end{dhkcode}

현재 \zsh{calculate}는 A의 wrapper 함수 객체를 참조하므로 A의 wrapper가 먼저 호출된다.

\begin{dhkzsh}
enter: A
\end{dhkzsh}

A의 wrapper는 자신의 클로저에 저장된 \zsh{func}를 호출한다.

그 객체는 B의 wrapper다.

\begin{dhkzsh}
enter: B
\end{dhkzsh}

B의 wrapper 역시 자신의 클로저에 저장된 \zsh{func}를 호출한다.

이번에는 원본 \zsh{calculate} 함수 객체가 호출된다.

\begin{dhkzsh}
body: 3
\end{dhkzsh}

원본 함수는 정수 객체 6을 반환한다.

그 반환값이 B의 wrapper로 전달되고, B의 wrapper는 \zsh{"exit: B"}를 출력한 뒤 다시 6을 반환한다.

이후 A의 wrapper가 \zsh{"exit: A"}를 출력하고 같은 반환 객체를 최종 호출자에게 전달한다.

따라서 최종 출력은 다음과 같다.

\begin{dhkzsh}
result: 6
\end{dhkzsh}

여기서 중요한 점은 적용 순서가 B → A이기 때문에 최종 호출 순서가 A → B가 된다는 것이다.

마지막에 적용된 A가 가장 바깥쪽 함수 객체가 되기 때문이다.

\subsubsection{함수 객체의 생성과 이름 바인딩을 추적하기}

이번에는 실행 결과가 아니라 객체 관계 자체를 분석해 보자.

앞에서 작성한 코드를 개념적으로 풀어 쓰면 다음과 같다.

\begin{dhkcode}
decorator_A = layer("A")
decorator_B = layer("B")


def calculate(x):
    return x * 2


original = calculate

inner = decorator_B(original)
outer = decorator_A(inner)

calculate = outer
\end{dhkcode}

이 코드는 핵심적인 객체 변환 관계를 나타낸 것이다.

실제 데코레이터 구문의 모든 평가·바인딩 세부사항을 그대로 재현한 것은 아니다.

이제 각 이름이 참조하는 객체를 정리해 보자.

\begin{dhktable}
|이름	|참조 객체|	역할|
|:---|:---|:---|
|decorator_A|	A 설정을 유지하는 데코레이터 함수|	바깥쪽 wrapper 생성|
|decorator_B|	B 설정을 유지하는 데코레이터 함수|	안쪽 wrapper 생성|
|original|	원본 함수 객체|	실제 계산 수행|
|inner|	B의 wrapper 함수 객체|	원본 함수 호출 중개|
|outer|	A의 wrapper 함수 객체|	B의 wrapper 호출 중개|
|calculate|	outer와 동일한 객체|	최종 공개 이름|
\end{dhktable}

이 구조에서 새로운 함수 객체는 여러 개 생성되지만 원본 함수 객체가 수정되는 것은 아니다.

데코레이터가 각각 새로운 wrapper를 생성하고 반환하며, 최종 이름이 가장 바깥쪽 객체에 바인딩되는 것이다.

따라서 다음과 같은 관계가 성립한다.

\begin{dhkcode}
calculate is outer  # True
outer is inner       # False
inner is original    # False
\end{dhkcode}

데코레이터가 wrapper 대신 원본 함수를 그대로 반환한다면 이러한 객체 수와 동일성 관계는 달라질 수 있다.

즉, 데코레이터를 여러 번 적용했다는 사실만으로 반드시 그 수만큼 새로운 함수 객체가 생성된다고 단정할 수는 없다.

\subsubsection{\zsh{__wrapped__} 체인을 통해 중첩 구조 확인하기}

\zsh{A.8}에서 학습한 \zsh{functools.wraps}를 사용하면 원본 함수 참조가 체인 형태로 유지된다.

앞의 \zsh{layer}와 \zsh{calculate}를 그대로 사용한다고 하자.

\begin{dhkcode}
import inspect

outer = calculate
inner = outer.__wrapped__
original = inner.__wrapped__

print(outer is inner)
print(inner is original)

print(inspect.unwrap(calculate) is original)

print(outer.__name__)
print(inner.__name__)
print(original.__name__)
\end{dhkcode}

예상 출력:

\begin{dhkzsh}
False
False
True
calculate
calculate
calculate
\end{dhkzsh}

여기서 흥미로운 점이 있다.

세 함수 객체는 모두 다른 객체인데 \zsh{__name__}은 같다.

왜냐하면 각 wrapper에 \zsh{@wraps(func)}가 적용되어 원본 쪽 함수의 메타데이터가 설정되었기 때문이다.

따라서 \zsh{__name__}만으로는 어떤 wrapper인지 구별할 수 없다.

반면 \zsh{__wrapped__}는 각 wrapper가 어떤 객체로 감싸도록 설정되었는지 추적할 수 있게 한다.

\begin{dhkgraph}[uniform]
node n1 (1,1) :: outer = A의 wrapper // calculate 이름이 참조하는 객체
node n2 (2,1) :: inner = B의 wrapper // A의 wrapper가 감싸는 객체
node n3 (3,1) :: original = 원본 calculate // 실제 x*2 계산을 수행하는 함수 객체
edge n1 -> n2 :: __wrapped__
edge n2 -> n3 :: __wrapped__

caption :: __wrapped__는 각 wrapper가 감싸는 객체를 가리킨다. 실제 호출도 이 경로를 따르는지는 wrapper의 본문에 달려 있다.
\end{dhkgraph}

\begin{dhknote}
\begin{dhksmry}[centered]
inspect.signature()의 동작
\end{dhksmry}
\end{dhknote}

다음 코드를 살펴보자.

\begin{dhkcode}
print(inspect.signature(calculate))

print(
    inspect.signature(
        calculate,
        follow_wrapped=False
    )
)
\end{dhkcode}

예상 출력:

\begin{dhkzsh}
(x)
(*args, **kwargs)
\end{dhkzsh}

첫 번째 결과는 \zsh{__wrapped__} 체인을 따라 원본 함수의 시그니처를 조회한 결과다.

두 번째 결과는 바깥쪽 wrapper 자체의 시그니처를 조회한 결과이다.

따라서 여러 데코레이터가 중첩되어 있어도 \zsh{functools.wraps}를 적절히 사용하면 외부 도구가 원본 함수의 인터페이스를 분석할 수 있다.

다만 중간 데코레이터가 \zsh{__wrapped__}를 설정하지 않으면 체인이 끊길 수 있다.

그리고 \zsh{__wrapped__} 체인은 메타데이터에 관한 관계이므로, wrapper가 실제로 \zsh{func()}를 호출하는지까지 보장하지는 않는다.

\subsubsection{\zsh{__closure__}로 실제 함수 참조 관계 검증하기}

이번에는 \zsh{A.4}와 \zsh{A.5}에서 학습한 클로저 셀을 활용해 본다.

앞의 \zsh{layer} 구현에서 wrapper는 다음 두 자유 변수를 참조한다.

\begin{dhkcode}
label
func
\end{dhkcode}

따라서 다음과 같이 클로저 셀을 확인할 수 있다.

\begin{dhkcode}
def closure_map(func):
    return dict(
        zip(
            func.__code__.co_freevars,
            func.__closure__ or ()
        )
    )


outer_cells = closure_map(calculate)
inner_cells = closure_map(calculate.__wrapped__)

print(outer_cells["label"].cell_contents)
print(inner_cells["label"].cell_contents)

print(
    outer_cells["func"].cell_contents
    is calculate.__wrapped__
)

print(
    inner_cells["func"].cell_contents
    is calculate.__wrapped__.__wrapped__
)
\end{dhkcode}

예상 출력:

\begin{dhkzsh}
A
B
True
True
\end{dhkzsh}

이 결과를 해석해 보자.

A의 wrapper에 저장된 \zsh{func}는 B의 wrapper 함수 객체를 참조한다.

반면 B의 wrapper에 저장된 \zsh{func}는 원본 \zsh{calculate} 함수 객체를 참조한다.

따라서 두 wrapper는 같은 이름 \zsh{func}를 자유 변수로 사용하지만 서로 다른 클로저 셀에 접근한다.

\begin{dhktable}
|함수 객체|	label 셀의 내용|	func 셀의 내용|
|:---|:---|:---|
|A의 wrapper|	"A"|	B의 wrapper|
|B의 wrapper|	"B"|	원본 calculate|
\end{dhktable}

여기서 클로저의 핵심 원리가 다시 나타난다.

자유 변수의 이름이 같다는 사실은 동일한 바인딩을 공유한다는 뜻이 아니다.

각 wrapper가 어느 데코레이터 호출에서 생성되었는지에 따라 클로저 셀이 달라진다.

\begin{dhknote}
\begin{dhksmry}[centered]
한 가지 더: 서로 다른 함수 객체가 같은 코드 객체를 사용할 수 있다
\end{dhksmry}
\end{dhknote}

앞의 \zsh{layer} 구현에서는 A와 B의 wrapper가 동일한 \zsh{def wrapper} 정의문에서 생성된다.

따라서 다음 결과를 확인할 수 있다.

\begin{dhkcode}
outer = calculate
inner = calculate.__wrapped__

print(outer is inner)
print(outer.__code__ is inner.__code__)
print(outer.__closure__ is inner.__closure__)
\end{dhkcode}

예상 출력:

\begin{dhkzsh}
False
True
False
\end{dhkzsh}

이 결과는 중요하다.

A의 wrapper와 B의 wrapper는 서로 다른 함수 객체다.

하지만 같은 함수 정의문의 코드 객체를 공유할 수 있다.

반면 클로저 환경은 서로 다르다.

즉, Python의 함수 객체를 분석할 때 다음 세 요소를 구별해야 한다.

\begin{dhksmry}
    - 함수 객체의 동일성

    - 실행할 코드 객체의 동일성

    - 함수 객체에 연결된 클로저 환경의 동일성
\end{dhksmry}

같은 코드 객체를 실행하더라도 서로 다른 클로저를 참조하면 동작이 달라질 수 있다.

이것이 함수 객체와 코드 객체를 구별해야 하는 이유 중 하나다.

\subsubsection{호출 프레임과 반환 순서}

이제 실행 환경의 관점에서 중첩 호출을 분석해 보자.

다음 호출을 생각해 보자.

\begin{dhkcode}
calculate(3)
\end{dhkcode}

실제 호출 경로는 다음과 같다.

\begin{dhkbranch}[in=8]
node main :: 호출자
arrow main
node main :: A의 wrapper 호출
line main
line main :: func(3)
arrow main
node main :: B의 wrapper 호출
line main
line main :: func(3)
arrow main
node main :: 원본 calculate 호출
line main
line main :: return 6
arrow main
node main :: B의 wrapper
line main
line main :: result -> 6
line main :: return result
arrow main
node main :: A의 wrapper
line main
line main :: result -> 6
line main :: return result
arrow main
node main :: 호출자
\end{dhkbranch}

여기서 wrapper마다 별도의 함수 호출이 발생하므로 각 호출은 자신의 지역 실행 환경을 가진다.

예를 들어 A의 wrapper에도 \zsh{args, kwargs, result}라는 지역 이름이 존재하고, B의 wrapper에도 같은 이름이 존재한다.

하지만 이 이름들이 하나의 네임스페이스를 공유하는 것은 아니다.

A의 \zsh{result}와 B의 \zsh{result}는 서로 다른 함수 호출의 지역 바인딩이다.

두 이름이 같은 정수 객체 6을 참조할 수 있을 뿐이다.

따라서 다음 원리가 성립한다.

\begin{dhksmry}[centered]
중첩된 wrapper가 같은 반환 객체를 전달하더라도

각 wrapper의 지역 이름은 독립적이다.
\end{dhksmry}

그리고 \zsh{return result}는 객체를 자동으로 복사하지 않는다.

각 wrapper가 반환 객체를 변경하거나 교체하지 않는다면 원본 함수가 반환한 객체가 그대로 호출자에게 전달될 수 있다.

\subsubsection{데코레이터의 적용 순서를 바꾸면 결과도 달라질까?}

지금까지 살펴본 데코레이터는 원본 함수의 반환값을 그대로 전달했다.

그렇다면 wrapper가 반환값을 변경하는 경우에는 어떻게 될까?

다음 두 데코레이터를 정의해 보자.

\begin{dhkcode}
from functools import wraps

def add_one(func):

    @wraps(func)
    def wrapper(x):
        return func(x) + 1

    return wrapper


def double(func):

    @wraps(func)
    def wrapper(x):
        return func(x) * 2

    return wrapper
\end{dhkcode}

이제 적용 순서를 다르게 해 보자.

\begin{dhkcode}
@double
@add_one
def first(x):
    return x


@add_one
@double
def second(x):
    return x


print(first(3))
print(second(3))
\end{dhkcode}

예상 출력:

\begin{dhkzsh}
8
7
\end{dhkzsh}

왜 서로 다른 결과가 발생할까?

\begin{dhknote}
\begin{dhksmry}[centered]
첫 번째 함수
\end{dhksmry}
\end{dhknote}

\begin{dhkcode}
@double
@add_one
def first(x):
    return x
\end{dhkcode}

이 코드는 다음과 대응한다.

\begin{dhkcode}
first = double(add_one(first))
\end{dhkcode}

실제 호출은 다음 순서로 진행된다.

\begin{dhkbranch}[in=8]
node main :: first(3)
arrow main
node main :: double의 wrapper
arrow main
node main :: add_one의 wrapper
arrow main
node main :: 원본 first(3)
handoff main -> a :: 3 반환
arrow a
node a :: add_one: 3 + 1 = 4
arrow a
node a :: double: 4 * 2 = 8
\end{dhkbranch}

따라서 8이 반환된다.

\begin{dhknote}
\begin{dhksmry}[centered]
두 번째 함수
\end{dhksmry}
\end{dhknote}

\begin{dhkcode}
@add_one
@double
def second(x):
    return x
\end{dhkcode}

이 코드는 다음과 대응한다.

\begin{dhkcode}
second = add_one(double(second))
\end{dhkcode}

실제 호출은 다음과 같다.

\begin{dhkbranch}[in=8]
node main :: second(3)
arrow main
node main :: add_one의 wrapper
arrow main
node main :: double의 wrapper 
arrow main
node main :: 원본 second(3)
handoff main -> a :: 3 반환
arrow a
node a :: double: 3 * 2 = 6
arrow a
node a :: add_one: 6 + 1 = 7
\end{dhkbranch}

따라서 7이 반환된다.

이 차이는 수학적인 함수 합성(Function Composition)과 연결된다.

원본 함수를 \(f\), 반환값에 1을 더하는 변환을 \(P\), 두 배로 만드는 변환을 \(D\)라고 하자.

그러면 첫 번째 경우는 다음과 같다.

\[
D(P(f))(x)=2(f(x)+1)
\]

두 번째 경우는 다음과 같다.

\[
P(D(f))(x)=2f(x)+1
\]

일반적으로 두 결과는 같지 않다.

즉, 데코레이터의 적용 순서는 교환법칙을 만족하지 않는다.

서로 다른 데코레이터를 적용하는 순서를 바꾸면 호출 구조뿐 아니라 최종 반환값도 달라질 수 있다.

따라서 여러 데코레이터를 사용할 때는 각 데코레이터가 인수, 반환값, 상태, 예외 중 무엇을 변경하는지 확인해야 한다.

\subsubsection{예외가 발생하면 wrapper는 어떤 순서로 종료될까?}

\zsh{A.7}에서는 예외가 발생하면 원본 함수의 정상적인 반환 경로가 중단된다는 사실을 학습했다.

여러 wrapper가 중첩된 경우에도 같은 규칙이 적용된다.

다음 코드를 살펴보자.

\begin{dhkcode}
from functools import wraps

def layer(label):

    def decorator(func):

        @wraps(func)
        def wrapper(*args, **kwargs):
            print(label, "enter")

            try:
                return func(*args, **kwargs)

            finally:
                print(label, "finally")

        return wrapper

    return decorator

@layer("A")
@layer("B")
def fail():
    print("body")
    raise ValueError("failed")

try:
    fail()

except ValueError:
    print("caught")
\end{dhkcode}

예상 출력:

\begin{dhkzsh}
A enter
B enter
body
B finally
A finally
caught
\end{dhkzsh}

이 실행 순서를 분석해 보자.

먼저 A의 wrapper가 호출된다.

A의 wrapper는 \zsh{"A enter"}를 출력하고 B의 wrapper를 호출한다.

B의 wrapper는 \zsh{"B enter"}를 출력하고 원본 \zsh{fail} 함수를 호출한다.

원본 함수는 \zsh{"body"}를 출력한 다음 \zsh{ValueError}를 발생시킨다.

여기서부터는 정상적인 반환이 아니라 예외 전파가 시작된다.

B의 wrapper는 \zsh{try} 블록에서 예외를 받지만 해당 예외를 처리하는 \zsh{except} 블록이 없다.

다만 \zsh{finally}가 존재하므로 \zsh{"B finally"}를 출력한다.

그다음 예외가 A의 wrapper로 전파된다.

A의 wrapper 역시 \zsh{finally}를 실행하여 \zsh{"A finally"}를 출력한다.

마지막으로 호출자의 \zsh{except ValueError} 블록이 예외를 처리한다.

\begin{dhkgraph}[uniform]
node n1 (1,1) :: A enter
node n2 (2,1) :: B enter
node n3 (3,1) :: 원본: 예외 발생
node n4 (3,2) :: B finally
node n5 (2,2) :: A finally
node n6 (1,2) :: 호출자의 except

edge n1 -> n2
edge n2 -> n3
edge n3 -> n4
edge n4 -> n5
edge n5 -> n6

explain :: 예외 발생 시 실행 경로 // 진입은 안쪽으로, 예외 전파는 바깥쪽으로
\end{dhkgraph}

여기서 중요한 점은 다음이다.

\begin{dhksmry}[centered]
예외가 발생하면 wrapper의 정상적인 반환 처리는 생략될 수 있지만,

finally는 예외가 전파되는 과정에서 실행된다.
\end{dhksmry}

다만 \zsh{finally} 자체에서 새로운 예외가 발생하거나 \zsh{return} 등으로 제어 흐름을 변경하면 기존 예외의 전파 방식도 달라질 수 있다.

따라서 여러 데코레이터의 예외 처리 구조를 분석할 때는 각 wrapper가 예외를 처리하는지, 다시 발생시키는지, 다른 예외로 변환하는지 구별해야 한다.

\subsubsection{실전 응용 - 로깅과 재시도의 적용 순서}

이제 \zsh{A.9}에서 학습한 재시도 데코레이터와 로깅 데코레이터를 결합해 보자.

이 예제는 데코레이터 순서가 단순한 문법 문제가 아니라 실제 프로그램의 관찰 가능한 동작을 결정한다는 사실을 보여준다.

\begin{dhkcode}
from functools import wraps

def log(func):

    @wraps(func)
    def wrapper(*args, **kwargs):
        print("LOG")
        return func(*args, **kwargs)

    return wrapper


def retry(max_attempts):

    def decorator(func):

        @wraps(func)
        def wrapper(*args, **kwargs):

            for attempt in range(max_attempts):

                try:
                    return func(*args, **kwargs)

                except ValueError:
                    if attempt == max_attempts - 1:
                        raise

        return wrapper

    return decorator
\end{dhkcode}

이제 같은 기능의 함수를 두 가지 순서로 데코레이션 해보자.

\begin{dhkcode}
state_a = {"count": 0}


@log
@retry(3)
def task_a():
    state_a["count"] += 1

    if state_a["count"] < 3:
        raise ValueError

    return "Success"


state_b = {"count": 0}


@retry(3)
@log
def task_b():
    state_b["count"] += 1

    if state_b["count"] < 3:
        raise ValueError

    return "Success"


print("A:", task_a())
print("B:", task_b())
\end{dhkcode}

예상 출력:

\begin{dhkzsh}
LOG
A: Success
LOG
LOG
LOG
B: Success
\end{dhkzsh}

왜 A에서는 로그가 한 번만 출력되고 B에서는 세 번 출력될까?

\begin{dhknote}
\begin{dhksmry}[centered]
사례 A - 로깅이 바깥쪽
\end{dhksmry}
\end{dhknote}

\begin{dhkcode}
@log
@retry(3)
def task_a():
    ...
\end{dhkcode}

호출 구조는 다음과 같다.

\begin{dhkdir}[in=6]
log wrapper
> retry wrapper
>> 원본 호출 1 실패
>> 원본 호출 2 실패
>> 원본 호출 3 성공
\end{dhkdir}

로깅 wrapper는 한 번 호출된다.

그 내부에서 재시도 wrapper가 세 번의 원본 함수 호출을 수행하는 것이다.

따라서 \zsh{"LOG"}는 한 번만 출력된다.

\begin{dhknote}
\begin{dhksmry}[centered]
사례 B - 재시도가 바깥쪽
\end{dhksmry}
\end{dhknote}

\begin{dhkcode}
@retry(3)
@log
def task_b():
    ...
\end{dhkcode}

호출 구조는 다음과 같다.

\begin{dhkdir}[in=6]
retry wrapper
> log wrapper -> 원본 호출 1 실패
> log wrapper -> 원본 호출 2 실패
> log wrapper -> 원본 호출 3 성공
\end{dhkdir}

이번에는 재시도 wrapper가 로깅 wrapper를 세 번 호출한다.

따라서 \zsh{"LOG"}도 세 번 출력된다.

이 차이는 실전에서 중요하다.

\begin{dhktable}
|목적	|적합한 구조|
|:---|:---|
|사용자 요청 한 건당 로그 한 번|	로깅을 재시도 바깥쪽에 배치|
|재시도 시도마다 로그 기록|	로깅을 재시도 안쪽에 배치|
\end{dhktable}

물론 실제 로깅 구현에서 무엇을 기록하는지에 따라 세부 설계는 달라질 수 있다.

핵심은 데코레이터를 어느 위치에 배치하는지가 기능의 적용 범위를 결정한다는 사실이다.

이 원리는 로깅뿐 아니라 캐싱, 인증, 트랜잭션, 성능 측정 등에도 적용된다.

\subsubsection{모든 데코레이터가 wrapper를 생성하는 것은 아니다}

지금까지는 모든 데코레이터가 새로운 wrapper를 반환하는 구조를 살펴봤다.

하지만 \zsh{A.6}에서 학습했듯 데코레이터는 원본 함수 객체를 그대로 반환할 수도 있다.

대표적인 사례가 등록형 데코레이터(Registration Decorator)다.

다음 코드를 살펴보자.

\begin{dhkcode}
from functools import wraps

registry = {}

def register(func):
    registry[func.__name__] = func
    return func


def trace(func):

    @wraps(func)
    def wrapper(*args, **kwargs):
        print("TRACE")
        return func(*args, **kwargs)

    return wrapper
\end{dhkcode}

이제 적용 순서를 달리해 보자.

\begin{dhkcode}
@register
@trace
def first():
    return 10


@trace
@register
def second():
    return 20


print(registry["first"] is first)

print(registry["second"] is second)

print(
    registry["second"]
    is second.__wrapped__
)
\end{dhkcode}

예상 출력:

\begin{dhkzsh}
True
False
True
\end{dhkzsh}

왜 이런 차이가 발생할까?

첫 번째 함수는 다음과 같이 반환된다.

\begin{dhkcode}
first = register(trace(first))
\end{dhkcode}

\zsh{trace(first)}가 먼저 실행되어 wrapper 함수 객체를 반환한다.

그다음 \zsh{register}가 그 wrapper 객체를 전달받아 등록한다.

\zsh{register}는 전달받은 객체를 그대로 반환하므로 최종 이름도 \zsh{first}도 등록된 wrapper를 참조한다.

따라서 다음 비교는 \zsh{True}다.

\begin{dhkcode}
registry["first"] is first
\end{dhkcode}

반면 두 번째 함수는 다음과 같다.

\begin{dhkcode}
second = trace(register(second))
\end{dhkcode}

이번에는 \zsh{register}가 원본 함수 객체를 먼저 등록한다.

그다음 \zsh{trace}가 새로운 wrapper를 생성해 반환한다.

따라서 레지스트리에는 원본 함수 객체가 저장되어 있지만 최종 이름 \zsh{second}는 wrapper를 참조한다.

그래서 다음 비교가 \zsh{False}인 것이다.

\begin{dhkcode}
registry["second"] is second
\end{dhkcode}

이 예제는 중요한 사실을 보여준다.

데코레이터의 적용 순서는 어떤 객체가 외부 시스템에 전달되거나 등록되는지도 결정할 수 있다.

이 내용은 \zsh{A.12}에서 FastAPI와 등록형 데코레이터를 분석할 때 다시 연결할 것이다.

\subsubsection{데코레이터 적용 중 예외가 발생하면 이름 바인딩은 어떻게 될까?}

이제 데코레이터가 적용되는 도중 오류가 발생하는 경계 조건을 살펴본다.

\begin{dhkcode}
def greet():
    return "Old"


def reject(func):
    print("reject called")
    raise ValueError("Decoration failed")


try:
    @reject
    def greet():
        return "New"

except ValueError:
    print("caught")


print(greet())
\end{dhkcode}

예상 출력:

\begin{dhkzsh}
reject called
caught
Old
\end{dhkzsh}

왜 \zsh{"New"}가 아니라 \zsh{"Old"}가 출력될까?

두 번째 \zsh{def greet()} 문장이 실행될 때 새로운 원본 함수 객체가 생성된다.

그 객체가 \zsh{reject}에 전달된다.

하지만 \zsh{reject}가 \zsh{ValueError}를 발생시키므로 데코레이터 적용이 정상적으로 완료되지 않는다.

따라서 새로운 반환 객체를 이름 \zsh{greet}에 바인딩하는 단계에 도달하지 못한다.

기존 모듈 이름 \zsh{greet}은 여전히 이전 함수 객체를 참조한다.

즉, 함수 정의문이 실행되는 과정에서 데코레이터 적용이 실패하면 새로운 함수 객체의 최종 이름 바인딩이 이루어지지 않을 수 있다.

이것은 데코레이터가 기존 함수를 제자리에서 수정하는 기능이 아니라, 최종 반환 객체를 이름에 바인딩하는 구조라는 사실을 다시 보여준다.

단, 데코레이터가 예외를 발생시키기 전에 수행한 다른 부작용까지 자동으로 취소되는 것은 아니다.

\subsubsection{데코레이터 평가·적용·호출을 구별하는 판단 기준}

지금까지의 내용을 하나의 표로 정리해 본다.

\begin{dhktable}
|질문|	확인할 대상|	판단 기준|
|:---|:---|:---|
|어떤 데코레이터 표현식이 먼저 실행되는가?|	@A(), @B()|	위에서 아래로 평가|
|어떤 데코레이터가 원본 함수를 먼저 전달받는가?|	A(B(func))|	아래쪽 B부터 적용|
|최종 이름은 어떤 객체를 참조하는가?|	가장 바깥쪽 적용 결과|	마지막 데코레이터의 반환 객체|
|어떤 wrapper가 먼저 호출되는가?|	최종 이름이 참조하는 객체|	바깥쪽 wrapper부터|
|원본 함수는 언제 호출되는가?|	각 wrapper의 실제 구현|	안쪽 wrapper가 원본을 호출할 때|
|반환값은 어떤 순서로 전달되는가?|	중첩된 호출의 반환|	안쪽에서 바깥쪽으로|
|예외는 어떤 순서로 전파되는가?|	try, except, finally|	호출 스택을 따라 바깥쪽으로|
|원본 함수 객체는 어디에 있는가?|	__wrapped__, 클로저 등|	실제 데코레이터 구현에 따라 다름|
|적용 순서를 바꿔도 결과가 같은가?|	인수·반환·상태·예외 처리|	일반적으로 보장되지 않음|
\end{dhktable}

여기서 특히 주의해야 할 점은 데코레이터 적용 순서와 wrapper 호출 순서를 같은 의미로 사용하지 않는 것이다.

데코레이터를 적용하는 과정에서는 함수 객체가 바깥쪽으로 쌓이고, 실제 호출에서는 그 바깥쪽부터 실행이 시작된다.

\subsubsection{기존 Python 객체 모델과의 연결}

이번 \zsh{A.10}은 데코레이터만의 독립적인 규칙을 학습한 것이 아니다.

지금까지 학습한 Python 객체 모델을 하나의 구조로 결합한 학습이다.

\begin{dhknote}
\begin{dhksmry}[centered]
함수 객체와 이름 바인딩 - A.1
\end{dhksmry}
\end{dhknote}

데코레이터가 반환한 객체가 최종 함수 이름에 바인딩된다.

이때 원본 함수 객체와 wrapper 함수 객체의 동일성은 별개다.

\begin{dhknote}
\begin{dhksmry}[centered]
렉시컬 스코프와 LEGB - A.2, A.3
\end{dhksmry}
\end{dhknote}

각 wrapper는 자신이 정의된 렉시컬 환경을 기준으로 외부 이름을 탐색한다.

호출자가 어떤 위치에서 함수를 실행하는지는 클로저의 자유 변수 탐색 경로를 변경하지 않는다.

\begin{dhknote}
\begin{dhksmry}[centered]
클로저와 셀 객체 - A.4, A.5
\end{dhksmry}
\end{dhknote}

각 wrapper는 외부 함수 객체나 설정값을 클로저 셀을 통해 유지할 수 있다.

같은 코드 객체를 공유하는 wrapper라도 서로 다른 클로저를 사용하면 다른 함수 객체를 호출할 수 있다.

\begin{dhknote}
\begin{dhksmry}[centered]
데코레이터 변환식 - A.6
\end{dhksmry}
\end{dhknote}

여러 데코레이터의 적용은 중첩된 함수 호출 형태로 표현할 수 있다.

\begin{dhkcode}
func = A(B(func))
\end{dhkcode}

\begin{dhknote}
\begin{dhksmry}[centered]
Wrapper의 인수 전달과 반환 - A.7
\end{dhksmry}
\end{dhknote}

각 wrapper는 독립적으로 인수를 전달하고 반환값을 처리한다.

따라서 여러 wrapper를 거치면서 반환값이 변경되거나 예외가 처리될 수 있다.

\begin{dhknote}
\begin{dhksmry}[centered]
functools.wraps와 메타데이터 - A.8
\end{dhksmry}
\end{dhknote}

\zsh{__wrapped__} 체인을 이용해 원본 함수 객체를 추적할 수 있다.

하지만 이 체인은 실제 호출을 강제하는 장치가 아니다.

\begin{dhknote}
\begin{dhksmry}[centered]
데코레이터 팩토리 - A.9
\end{dhksmry}
\end{dhknote}

\zsh{@A()}와 \zsh{@B()}는 먼저 각각 평가되어 데코레이터 객체를 준비한다.

그다음 아래쪽부터 원본 함수 객체에 적용된다.

결국 \zsh{A.10}은 이러한 개념들을 하나의 실행 흐름으로 연결하는 것이다.

\subsubsection{핵심 정리}

이번 학습에서 반드시 구별해야 하는 것은 네 가지다.

\begin{dhknote}
\begin{dhksmry}[centered]
첫째, 데코레이터 표현식의 평가 순서
\end{dhksmry}
\end{dhknote}

\begin{dhkcode}
@A()
@B()
def func():
    ...
\end{dhkcode}

\zsh{A()}가 먼저 평가되고 그다음 \zsh{B()}가 평가된다.

\begin{dhknote}
\begin{dhksmry}[centered]
둘째, 데코레이터의 적용 순서
\end{dhksmry}
\end{dhknote}

원본 함수 객체가 생성된 다음 B가 먼저 적용되고 A가 나중에 적용된다.

\begin{dhkcode}
func = A_decorator(B_decorator(original))
\end{dhkcode}

\begin{dhknote}
\begin{dhksmry}[centered]
셋째, 실제 wrapper 호출 순서
\end{dhksmry}
\end{dhknote}

최종 이름이 A의 wrapper를 참조한다면 A의 wrapper가 먼저 호출된다.

\begin{dhksmry}[centered]
A wrapper \ra B wrapper \ra 원본 함수
\end{dhksmry}

\begin{dhknote}
\begin{dhksmry}[centered]
넷째, 정상 반환과 예외 전파 순서
\end{dhksmry}
\end{dhknote}

정상적인 반환은 원본 함수에서 시작해 B의 wrapper, A의 wrapper 순서로 전달된다.

예외 역시 처리되지 않는다면 호출 스택을 따라 바깥쪽으로 전파된다.

이 네 가지를 이해하면 여러 데코레이터가 적용된 코드를 볼 때 단순히 출력 순서를 암기하는 것이 아니라, 각 함수 객체가 어떤 객체를 참조하고 어떤 시점에 호출되는지를 근거로 실행 과정을 예측할 수 있다.
